\documentclass[10pt,a4paper]{amsart}
\usepackage[margin=19mm]{geometry}
\usepackage{amsmath,amssymb,mathtools}
\usepackage[scr=rsfs,cal=euler]{mathalfa}
\usepackage{xfrac}
\usepackage[unicode,hidelinks,bookmarks=false]{hyperref}
\newcommand{\ie}{i.e.\ }
\newcommand{\cf}{cf.\ }
\newcommand{\resp}{respectively\ }
\newcommand{\Subsection}{\subsection}
\providecommand{\nobreakdash}{\nobreak\hbox{-}}
\setlength{\emergencystretch}{2em}
\multlinegap0pt
\usepackage{amssymb}
\usepackage{enumerate}
\usepackage{xcolor}
\usepackage{tikz}
\newtheorem{lemma}{Lemma}
\newcommand{\R}{{ \mathbb{R}  }}
\newcommand{\bke}[1]{\left( #1 \right)}
\newcommand{\bkt}[1]{\left[ #1 \right]}
\newcommand{\norm}[1]{\left\Vert #1 \right\Vert}
\title{Keller--Segel--Navier--Stokes systems involving general sensitivities with Signal-Dependent Power-Law Decay}
\author{Jaewook Ahn, Sukjung Hwang}
\date{Source: arXiv:2603.07870v1 (CC BY 4.0)}
\begin{document}
\maketitle

\noindent\emph{Source and licence.} Keller--Segel--Navier--Stokes systems involving general sensitivities with Signal-Dependent Power-Law Decay. Version arXiv:2603.07870v1 (CC BY 4.0), distributed by arXiv under the Creative Commons Attribution 4.0 International licence, http://creativecommons.org/licenses/by/4.0/, as recorded in the arXiv licence metadata for this exact version and shown on the abstract page https://arxiv.org/abs/2603.07870v1. Full licence notice: \texttt{licenses/arxiv-2603.07870-CC-BY-4.0.txt}.\par\smallskip

\noindent\textit{Editorial scope.} Counted unit: the article's lemma LEMLOCINT (source0.tex lines 913--932). The article's complete original proof of it is delivered with it (source0.tex lines 935--1076, one \texttt{proof} environment of 142 lines). No mathematics is written, repaired or re-derived: the statement and the proof are the article's own lines, in the article's own notation, and no retained line is substituted or re-flowed.  Retained uncounted as the source-local context the proof needs: lines 313--321; lines 448--450; lines 817--833; lines 841--843; lines 863--865; lines 871--873.  Nothing was omitted from any retained span, so the package carries no omission record.  No retained line contains a citation, so the wrapper carries no bibliography and no bibliography entry.  No retained line contains a figure environment, an image-inclusion command or drawing code, so no image is delivered and the package declares no asset dependency.  Editorial formatting only, in addition: an AMS \texttt{amsart} preamble that restates the article's own declarations and macros used by the retained lines, namely the environments \texttt{lemma} on separate counters, together with the standard packages those lines need.  The retained environments print in this document as Lemma 1 in order of appearance, whereas the article prints the same items with the numbering of its own full document.  The complete native source of the article as distributed in the arXiv e-print for this exact version is bundled unchanged as \texttt{sources/195976/source0.tex}.
\begin{equation}\label{KS01}
\begin{cases}
n_t + u\cdot\nabla n = \Delta n - \nabla\cdot\bigl(n\,S(x,n,c)\nabla c\bigr), & \text{in } \Omega\times(0,\infty),\\
-\Delta c + u\cdot\nabla c + c = n, & \text{in } \Omega\times(0,\infty),\\
u_t + (u\cdot\nabla)u = \Delta u - \nabla\pi + n\nabla\Phi, & \text{in } \Omega\times(0,\infty),\\
\nabla\cdot u = 0, & \text{in } \Omega\times(0,\infty), \\
n(\cdot, 0)=n_0, \ u(\cdot, 0)=u_0, & \text{in } \Omega,
\end{cases}
\end{equation}

\begin{equation}\label{NABLACL2}
r\int_{\Omega}\frac{|\nabla c|^{2}}{(\sigma_{1}+c)^{r+1}}(\cdot,t)+\int_{\Omega}\frac{n}{(\sigma_{1}+c)^{r}}(\cdot,t)=\int_{\Omega}\frac{c}{(\sigma_{1}+c)^{r}}(\cdot,t),\qquad t<T_{\rm max}.
\end{equation}

\begin{lemma}\label{LEMSTEST}
Let $\delta >0$. There exists a radially decreasing function 
$\varphi_{\delta} \in C_{0}^{\infty}(\mathbb{R}^d)$, $d\ge1$,  such that 
	\[
	\varphi_{\delta}(x) = 
	\begin{cases}
		1, & \text{if } x\in B_{\frac{\delta}{2}}(0), \\
		0, & \text{if } x\in \mathbb{R}^d \setminus B_{\delta}(0),
	\end{cases}
	\]
	$0\leq \varphi_{\delta} \leq 1 $    in $\mathbb{R}^d$,
	and 
	\[
 |\nabla \varphi_{\delta}| \leq K \varphi_{\delta}^{\frac{1}{2}} \ \text{ in } \ \mathbb{R}^d,
	\]
	where $K$ is a positive constant of order $\mathcal{O}(\delta^{-1})$.
\end{lemma}

 \begin{equation}\label{CLOWER}
\inf_{\overline{\Omega}}\{s_{1}+c(\cdot,t)\}\ge C\quad\mbox{for all }t<T_{\rm max}.
 \end{equation}

\begin{equation}\label{FunctionF}
 F(c):=  \int_{1}^{c}\frac{1}{(s_{1}+s)^{k}}ds \quad  \mbox{for } k>\frac{1}{2}.
\end{equation}

\begin{equation}\label{FLEM}
 \|F(c)(\cdot,t)\|_{L^{p}(\Omega)}+\|\nabla F(c)(\cdot,t)\|_{L^{2}(\Omega)} \le C\quad\mbox{for all }t<T_{\rm max},
\end{equation}

\begin{lemma}\label{LEMLOCINT}
Let $k>\frac{1}{2}$.
Assume that either $s_{1}>0$, or $s_{1}\ge0$ and the system is fluid-free.
Denote $\varphi(x)=\varphi_{\delta}(x-q)$ and $B_{\delta}=B_{\delta}(q)$  for $q\in\overline{\Omega}$, where $\varphi_{\delta}$, $\delta>0$,  is the function given in   Lemma~\ref{LEMSTEST}. 
Then, there exist two positive constants $C_{1}=C_1(\Omega, \int_{\Omega} n_0) >0$ and  $C_2=C_2(\Omega, k, s_1, \int_{\Omega} n_0)>0$ independent of $\delta$ and $q$  such that
\begin{equation}\label{LEMLOCINT1}
\int_{\Omega}n^{2}\varphi^{4}
 \le C_{1} \bke{  \int_{\Omega} \frac{|\nabla n|^{2}}{n}\varphi^{4}  +\| \varphi^{2}  \|_{W^{1,\infty}(\R^{2})}^{2}},
\end{equation}
\begin{align}\label{LEMLOCINT2}
\begin{aligned} 
 \int_{\Omega} \frac{|\nabla c|^{4}}{(s_{1}+c)^{4k}}\varphi^{4}
 &\le  C_{2}\norm{\frac{\nabla c}{(s_{1}+c)^{k}}}_{L^{2}(\Omega\cap B_{\delta})} \norm{\frac{\nabla c}{(s_{1}+c)^{\frac{3-k}{2}}}}_{L^{2}(\Omega\cap B_{\delta})}^{2}  \int_{\Omega} \frac{|\nabla c|^{4}}{(s_{1}+c)^{4k}}\varphi^{4}
   \\
 &\quad +  C_{2}\norm{\frac{\nabla c}{(s_{1}+c)^{k}}}_{L^{2}(\Omega\cap B_{\delta})}
\bke{   \int_{\Omega} \frac{|\nabla n|^{2}}{n}\varphi^{4}  +\| \varphi^{2}  \|_{W^{1,\infty}(\R^{2})}^{2}  +\|\varphi^{\frac{4}{3}}\|_{W^{2,2}(\R^{2})}^{3}} \\
&\quad +C_{2} \norm{\frac{\nabla c}{(s_{1}+c)^{k}}}^{4}_{L^{2}(\Omega\cap B_{\delta})} \bke{\int_{\Omega}|u|^{\frac{12}{5}} }^{5}
\end{aligned}
\end{align}
\end{lemma}

\begin{proof}
First, to prove \eqref{LEMLOCINT1}, we use the Sobolev embedding $W^{1, 1}(\Omega)\hookrightarrow L^{2}(\Omega)$ valid in two dimensions that there exists $C=C(\Omega)>0$ satisfying
\[
 \int_{\Omega}n^{2}\varphi^{4}   \le C\bke{ \|\nabla(n\varphi^{2})\|_{L^{1}(\Omega)}^{2}+\|n\varphi^{2}\|_{L^{1}(\Omega)}^{2}}.
\]
Using   H\"older's inequality and  $\int_{\Omega}n=\int_{\Omega}n_{0}$, there is $C>0$, independent of $\delta$ and $q$, such that 
\[
\begin{aligned}
 \|\nabla(n\varphi^{2})\|_{L^{1}(\Omega)}^{2} +\|n\varphi^{2}\|_{L^{1}(\Omega)}^{2}
\le C\bke{ \|n_{0}\|_{L^{1}(\Omega)}  \int_{\Omega} \frac{|\nabla n|^{2}}{n}\varphi^{4} +\|n_{0}\|_{L^{1}(\Omega)}^{2}\| \varphi^{2}  \|_{W^{1,\infty}(\R^{2})}^{2} }.
 \end{aligned}
\]
We prove \eqref{LEMLOCINT1} by combining two inequalities above.

Recall $F$ from \eqref{FunctionF}, so that $\nabla F(c)=\frac{\nabla c}{(s_1+c)^k}$.
Now, using  H\"older's  inequality and $(a_{1}+a_{2})^{3}\le 4(a_{1}^{3}+a_{2}^{3})$ for $a_{1},a_{2}\ge0$,  we compute the following 
\begin{equation}\label{LEM5_1}
\begin{aligned}[b]
 \int_{\Omega}\frac{|\nabla c|^{4}}{(s_{1}+c)^{4k}}\varphi^{4}
&\le   \norm{\frac{\nabla c}{(s_{1}+c)^{k}}}_{L^{2}(\Omega\cap B_{\delta})}\bke{\int_{\Omega}\frac{|\nabla c|^{6}}{(s_{1}+c)^{6k}}\varphi^{8}}^{\frac{1}{2}}\\
%&=  \norm{\frac{\nabla c}{(s_{1}+c)^{k}}}_{L^{2}(\Omega\cap B_{\delta})} \norm{ \nabla F(c) \varphi^{\frac{4}{3}} }_{L^{6}(\Omega)}^{3}\\
 &= \norm{\frac{\nabla c}{(s_{1}+c)^{k}}}_{L^{2}(\Omega\cap B_{\delta})}\norm{ \nabla (F(c)\varphi^{\frac{4}{3}})- F(c) \nabla \varphi^{\frac{4}{3}} }_{L^{6}(\Omega)}^{3}
\\
 &\le   \norm{\frac{\nabla c}{(s_{1}+c)^{k}}}_{L^{2}(\Omega\cap B_{\delta})}\bke{\norm{  F(c)  \varphi^{\frac{4}{3}}}_{W^{1,6}(\Omega)} + \norm{F(c)\nabla \varphi^{\frac{4}{3}}}_{L^{6}(\Omega)}}^{3}
\\
& \le  4\norm{\frac{\nabla c}{(s_{1}+c)^{k}}}_{L^{2}(\Omega\cap B_{\delta})} \bke{\norm{  F(c)  \varphi^{\frac{4}{3}}}_{W^{1,6}(\Omega)}^{3}+  \norm{ F(c)}_{L^{12}(\Omega)}^{3}\norm{ \nabla \varphi^{\frac{4}{3}}}_{L^{12}(\Omega)}^{3}}.
\end{aligned}
\end{equation}
Note that $\norm{ F(c)}_{L^{12}(\Omega)}$ is uniformly bounded by Lemma~\ref{FLEM} with $p=12$.

First, we estimate $\norm{   F(c)  \varphi^{\frac{4}{3}}}_{W^{1,6}(\Omega)}$.
The Sobolev embedding $W^{2,\frac{3}{2}}(\Omega)\hookrightarrow W^{1,6}(\Omega)$ along with 
%the standard elliptic regularity theory 
a standard $W^{2,p}$-estimate with $p=\frac{3}{2}$
yields that there exists $C=C(\Omega)>0$ independent of $\delta$ and $q$ satisfying
\begin{equation}\label{LEM5_2}
\begin{aligned}
 \norm{   F(c)  \varphi^{\frac{4}{3}}}_{W^{1,6}(\Omega)}^{3}
 \le C\bke{  \norm{\Delta \bke{  F(c)   \varphi^{\frac{4}{3}} } }_{L^{\frac{3}{2}}(\Omega)}^{3} +\norm{    F(c)   \varphi^{\frac{4}{3}}}_{L^{\frac{3}{2}}(\Omega)}^{3} +\norm{ F(c)\nabla \varphi^{\frac{4}{3}} \cdot \nu }_{W^{\frac{1}{3},\frac{3}{2}}(\partial\Omega)}^{3} },
\end{aligned}
\end{equation}
where we used
$\nabla(F(c)\varphi^{\frac{4}{3}})\cdot\nu=F(c)\nabla \varphi^{\frac{4}{3}} \cdot \nu$ on $\partial\Omega$ since $\partial_{\nu} c = 0$ on $\partial \Omega$. 
By H\"older's inequality and   $W^{1,2}(\Omega)\hookrightarrow L^{\frac{12}{7}}(\Omega)$, we can find $C=C(\Omega)>0$ such that
\[
\| F(c)   \varphi^{\frac{4}{3}}\|_{L^{\frac{3}{2}}(\Omega)}^{3}\le C \norm{ F(c)}_{L^{12}(\Omega)}^{3}\|\varphi^{\frac{4}{3}}\|_{W^{2,2}(\R^{2})}^{3}.
\]
Note   from the trace inequality and a direct computation that
\begin{align*}
\begin{aligned}
\| F(c)\nabla \varphi^{\frac{4}{3}} \cdot \nu \|_{W^{\frac{1}{3},\frac{3}{2}}(\partial\Omega)}^{3}\le C \|F(c)\nabla \varphi^{\frac{4}{3}} \|_{W^{1,\frac{3}{2}}( \Omega)}^{3}.
\end{aligned}
\end{align*}
Note also from H\"older's inequality and $W^{1,2}(\Omega)\hookrightarrow L^{q}(\Omega)$ for all $q\in[1,\infty)$ that with some $C=C(\Omega)>0$
\[
\|F(c)\nabla \varphi^{\frac{4}{3}} \|_{W^{1,\frac{3}{2}}( \Omega)}^{3}\le C(\|F(c)\|_{L^{12}(\Omega)}^{3}+\|\nabla F(c)\|_{L^{2}(\Omega)}^{3}+ \|F(c)\|_{L^{6}(\Omega)}^{3})\|\varphi^{\frac{4}{3}}\|_{W^{2,2}(\R^{2})}^{3}.
\]
By applying Lemma~\ref{FLEM} with $p=12$ and $p=6$, the last two terms of \eqref{LEM5_2} are bounded uniformly. 
Hence it remains to estimate $\|\Delta(F(c)\varphi^{4/3})\|_{L^{3/2}}^{3}$.
Using  H\"older's inequality,  $(a_{1}+a_{2})^{3}\le 4(a_{1}^{3}+a_{2}^{3})$ for $a_{1},a_{2}\ge0$,  $\eqref{KS01}_{2}$, and $W^{2,2}(\Omega)\hookrightarrow W^{1,6}(\Omega)$, we compute 
\begin{align*}
\begin{aligned}
\|  \Delta(F(c)   \varphi^{\frac{4}{3}})  \|_{L^{\frac{3}{2}}(\Omega)}^{3}
&\le 4\| \Delta F(c)   \varphi^{\frac{4}{3}}   \|_{L^{\frac{3}{2}}(\Omega)}^{3}
    +4 \left( 2\| \nabla F(c) \cdot \nabla \varphi^{\frac{4}{3}}\|_{L^{\frac{3}{2}}(\Omega)}+\| F( c)  \Delta \varphi^{\frac{4}{3}}   \|_{L^{\frac{3}{2}}(\Omega)}\right)^{3}
\\
&\le  16\bke{\norm{\frac{ u\cdot \nabla c-n+c  }{(s_{1}+c)^{k}}  \varphi^{\frac{4}{3}}   }_{L^{\frac{3}{2}}(\Omega)}^{3}
+  k^{3}\norm{ \frac{|\nabla c|^{2}}{(s_{1}+c)^{k+1}}      \varphi^{\frac{4}{3}}   }_{L^{\frac{3}{2}}(\Omega)}^{3}} 
\\&\quad+ 8\left( \|\nabla F(c)  \|_{L^{2}(\Omega )} +\|F(c)\|_{L^{6}(\Omega)} \right)^{3} \|\varphi^{\frac{4}{3}}\|^{3}_{W^{2,2}(\R^{2})}  .
\end{aligned}
\end{align*}
The last term is bounded uniformly by Lemma~\ref{FLEM} with $p=6$.
By applying H\"{o}lder inequality, we note that 
\begin{align*}
\begin{aligned}
\norm{ \frac{|\nabla c|^{2}}{(s_{1}+c)^{k+1}}      \varphi^{\frac{4}{3}}   }_{L^{\frac{3}{2}}(\Omega)}^{3}  \le \norm{\frac{\nabla c}{(s_{1}+c)^{\frac{3-k}{2}}}}_{L^{2}(\Omega\cap B_{\delta})}^{2}  \int_{\Omega} \frac{|\nabla c|^{4}}{(s_{1}+c)^{4k}}\varphi^{4}. 
\end{aligned}
\end{align*}
Now, it follows that, by the H\"{o}lder inequality along with $ \varphi \le 1$, 
\begin{align*}
\begin{aligned}
\norm{\frac{ u\cdot \nabla c  }{(s_{1}+c)^{k}}  \varphi^{\frac{4}{3}}   }_{L^{\frac{3}{2}}(\Omega)}^{3}
\le \bke{\int_{\Omega}|u|^{\frac{12}{5}} }^{\frac{5}{4}}\bke{\int_{\Omega} \frac{|\nabla c|^{4}}{(s_{1}+c)^{4k}}\varphi^{4}}^{\frac{3}{4}}.
\end{aligned}
\end{align*}
There is a constant $C=C(\Omega, s_1, k, \int_{\Omega}n_0)>0$, by H\"older's inequality along with \eqref{NABLACL2} and \eqref{CLOWER}, such that 
\begin{align*}
\begin{aligned}
\norm{ \frac{ n }{(s_{1}+c)^{k}}  \varphi^{\frac{4}{3}}   }_{L^{\frac{3}{2}}(\Omega)}^{3} 
&\le  \norm{\frac{n}{(s_{1}+c)^{2k}}}_{L^{1}( \Omega \cap B_{\delta} )}\int_{\Omega}n^{2}\varphi^{4} 
\\&  \le \left(\int_{\Omega} \frac{1}{(s_{1}+c)^{2k-1}} \right) \int_{\Omega}n^{2}\varphi^{4} 
\\&\le  C\int_{\Omega}n^{2}\varphi^{4},
\end{aligned}
\end{align*}
Similarly, by H\"older's inequality, \eqref{CLOWER}, $\int_{\Omega}c=\int_{\Omega}n_{0}$, and $W^{2,2}(\Omega)\hookrightarrow L^{6}(\Omega)$,
there is $C>0$ such that % {\color{dred} $C=C(\Omega, s_1, \int_{\Omega} n_0)>0$}
\begin{align*}
\begin{aligned}
\norm{ \frac{ c }{(s_{1}+c)^{k}}  \varphi^{\frac{4}{3}}   }_{L^{\frac{3}{2}}(\Omega)}^{3} 
&\le \norm{ \frac{ c }{(s_{1}+c)^{k}}}_{L^{2}(\Omega \cap B_{\delta} )}^{3} \norm{  \varphi^{\frac{4}{3}}   }_{L^{6}(\Omega)}^{3} 
\\&  \le \left(\int_{\Omega} \frac{c}{(s_1 + c)^{2k-1}}\right)^{\frac{3}{2}} \norm{  \varphi^{\frac{4}{3}}   }_{L^{6}(\Omega)}^{3} 
\\&\le C \norm{\varphi^{\frac{4}{3}}}_{W^{2,2}(\R^{2})}^{3} .
\end{aligned}
\end{align*}
Therefore, there is a constant $C>0$ such that 
\[\begin{aligned}
\| \Delta(F(c)   \varphi^{\frac{4}{3}})  \|_{L^{\frac{3}{2}}(\Omega)}^{3}
&\leq C\int_{\Omega}n^{2}\varphi^{4} 
+ 16 \bke{\int_{\Omega}|u|^{\frac{12}{5}} }^{\frac{5}{4}}\bke{\int_{\Omega} \frac{|\nabla c|^{4}}{(s_{1}+c)^{4k}}\varphi^{4}}^{\frac{3}{4}} \\
&\quad + 16 k^3 \norm{\frac{\nabla c}{(s_{1}+c)^{\frac{3-k}{2}}}}_{L^{2}(\Omega\cap B_{\delta})}^{2}  \int_{\Omega} \frac{|\nabla c|^{4}}{(s_{1}+c)^{4k}}\varphi^{4}
 + C \norm{\varphi^{\frac{4}{3}}}_{W^{2,2}(\R^{2})}^{3}.
\end{aligned}\]
Moreover, there is a constant $C>0$ such that 
\[\begin{aligned}
\norm{   F(c)  \varphi^{\frac{4}{3}}}_{W^{1,6}(\Omega)}^{3}
&\leq C\int_{\Omega}n^{2}\varphi^{4} 
+ C\bke{\int_{\Omega}|u|^{\frac{12}{5}} }^{\frac{5}{4}}\bke{\int_{\Omega} \frac{|\nabla c|^{4}}{(s_{1}+c)^{4k}}\varphi^{4}}^{\frac{3}{4}} \\
&\quad + C k^3 \norm{\frac{\nabla c}{(s_{1}+c)^{\frac{3-k}{2}}}}_{L^{2}(\Omega\cap B_{\delta})}^{2}  \int_{\Omega} \frac{|\nabla c|^{4}}{(s_{1}+c)^{4k}}\varphi^{4}
 + C \norm{\varphi^{\frac{4}{3}}}_{W^{2,2}(\R^{2})}^{3}.
\end{aligned}\]

Combining the above estimates
shows that there exists a constant $C>0$ independent of $\delta$ and $q$ such that
\begin{align*}
\begin{aligned}
 \int_{\Omega}\frac{|\nabla c|^{4}}{(s_{1}+c)^{4k}}\varphi^{4} 
 & \le  C\norm{\frac{\nabla c}{(s_{1}+c)^{k}}}_{L^{2}(\Omega\cap B_{\delta})} \norm{\frac{\nabla c}{(s_{1}+c)^{\frac{3-k}{2}}}}_{L^{2}(\Omega\cap B_{\delta})}^{2}  \int_{\Omega} \frac{|\nabla c|^{4}}{(s_{1}+c)^{4k}}\varphi^{4}
 \\ & \quad +
  C \norm{\frac{\nabla c}{(s_{1}+c)^{k}}}_{L^{2}(\Omega\cap B_{\delta})}   \bkt{ \bke{\int_{\Omega} \frac{|\nabla c|^{4}}{(s_{1}+c)^{4k}}\varphi^{4}}^{\frac{3}{4}}\bke{\int_{\Omega}|u|^{\frac{12}{5}} }^{\frac{5}{4}} + \int_{\Omega}n^{2}\varphi^{4}  +\|\varphi^{\frac{4}{3}}\|_{W^{2,2}(\R^{2})}^{3}}.
\end{aligned}
\end{align*}
Finally, by Young's inequality, \eqref{NABLACL2} and   \eqref{LEMLOCINT1}, it follows that there exists a constant $C=C(\Omega, s_1, k, \int_{\Omega}n_0)>0$ independent of $\delta$ and $q$ satisfying
\[
\begin{aligned}
  \int_{\Omega}\frac{|\nabla c|^{4}}{(s_{1}+c)^{4k}}\varphi^{4} 
 &  \le C\norm{\frac{\nabla c}{(s_{1}+c)^{k}}}_{L^{2}(\Omega\cap B_{\delta})} \norm{\frac{\nabla c}{(s_{1}+c)^{\frac{3-k}{2}}}}_{L^{2}(\Omega\cap B_{\delta})}^{2}  \int_{\Omega} \frac{|\nabla c|^{4}}{(s_{1}+c)^{4k}}\varphi^{4}
 \\&\quad  +   C \norm{\frac{\nabla c}{(s_{1}+c)^{k}}}_{L^{2}(\Omega\cap B_{\delta})}     \bke{  \int_{\Omega} \frac{|\nabla n|^{2}}{n}\varphi^{4}  +\| \varphi^{2}  \|_{W^{1,\infty}(\R^{2})}^{2}  +\|\varphi^{\frac{4}{3}}\|_{W^{2,2}(\R^{2})}^{3}}
 \\&\quad + C \norm{\frac{\nabla c}{(s_{1}+c)^{k}}}^{4}_{L^{2}(\Omega\cap B_{\delta})} \bke{\int_{\Omega}|u|^{\frac{12}{5}} }^{5},
 \end{aligned}
\]
namely, \eqref{LEMLOCINT2}.
\end{proof}

\end{document}
